\section{A Benchmark from Issues, Logs, and Fix Commits}
\label{sec:benchmark}

The challenge asks for evaluation on a benchmark that contains issue
descriptions, logs, source code, and ground-truth faulty code regions.
Existing resources supply some of these but not the combination we need.
Fault localization benchmarks such as Defects4J~\cite{just2014defects4j}
supply failing tests, which is exactly the execution evidence that is
missing when an issue is first reported. Bench4BL pairs reports with
code and gold files at scale but does not select for reports that carry
logs~\cite{lee2018bench4bl}, and benchmarks built from pull
requests~\cite{aleithan2024swebenchplus} inherit solution text in the
issue thread. Pathidea assembled the closest resource, a log-bearing
set of Java issues~\cite{chen2022pathidea}.

What we add is a set selected for log-bearing reports, spanning four
languages rather than one, stratified by whether the report names any
faulty file, and shipped with a per-project pool of past incidents and
their fix commits so that history-based methods can be evaluated under a
stated temporal firewall. This section describes how each instance is
constructed and what quality controls we applied.

\subsection{Data sources}
\label{sec:benchmark:sources}

We join two public datasets. The first is World of Logs v2 (WoL), a
dataset of software logs extracted from online documents~\cite{wang2026wol}.
For every public Jira issue it records whether the issue text contains
log messages, and an uncertainty score for that extraction. The second is
The Public Jira Dataset v7 (PJD)~\cite{montgomery2022jira}, which provides
the complete metadata of 2.7 million issues from 16 public Jira trackers,
including creation time, resolution time, resolution status, and typed
links between issues. WoL is itself built on PJD, so the two join exactly
on the issue key. Both datasets are licensed CC-BY-4.0, which allows us to
redistribute the benchmark with attribution.

From the 16 trackers we selected four projects: Apache Kafka (Java and
Scala), MariaDB server (C and C++), MongoDB server (C++), and Qt
(C++, the \texttt{qtbase} module). The four projects come from four
different organizations and two tracker platforms, and they differ
strongly in how their users report bugs, as \cref{sec:benchmark:strata}
shows.

\subsection{From an issue to its ground truth}
\label{sec:benchmark:linkage}

An issue enters the benchmark if it satisfies four conditions. First, its
text contains logs according to WoL, with extraction uncertainty at most
0.05, the threshold the WoL authors recommend. Second, it is a bug: its
issue type is \emph{Bug} and its resolution is \emph{Fixed} or
\emph{Done}. Third, it can be linked to the commit that fixed it. For
Kafka, MariaDB, and MongoDB, developers put the issue key in the commit
subject (for example, ``KAFKA-10089 The stale ssl engine factory is not
closed after reconfigure''). Qt developers instead record the key in a
\texttt{Fixes:} trailer in the commit message. We scan the default branch
of each project's repository and keep issues linked to one to three
commits. Issues linked to more commits are long-running efforts rather
than single faults, and we drop them. We also drop revert commits and
merge commits, whose diffs do not describe a single fix. Fourth, the fix
must change source code: at least one modified file must be a non-test
source file of the project's main languages.

For every surviving issue, the \emph{ground truth} (gold) is the set of
non-test source files that the fix commit modified. The \emph{code
snapshot} is the repository at the parent of the fix commit. This is the
code as it existed while the bug was still present, so the fix itself is
never visible to a localization method. The \emph{query} is the issue
summary and description, truncated at 20{,}000 characters. This matches
the challenge input: a natural-language report with the logs the reporter
pasted into it.

\begin{table}[t]
\caption{Benchmark datasheet. ``Eligible'' counts log-bearing fixed bugs
(conditions 1 and 2); ``linked'' additionally requires a clean fix-commit
link (condition 3); ``final'' requires modified non-test source files
(condition 4). ``Hard'' instances do not name any gold file in their
text. The memory pool (\cref{sec:benchmark:memory}) does not require
logs, so it is larger than the benchmark itself.}
\label{tab:datasheet}
\centering
\begin{tabular}{lrrrr}
\toprule
 & Kafka & MariaDB & MongoDB & Qt \\
\midrule
Eligible issues        & 1{,}172 & 3{,}004 & 6{,}407 & 7{,}412 \\
Linked (1--3 commits)  & 827     & 2{,}010 & 5{,}093 & 713 \\
Final instances        & 680     & 1{,}683 & 3{,}618 & 456 \\
\midrule
Hard stratum           & 28\%    & 44\%    & 60\%    & 68\% \\
With stack trace       & 47\%    & 57\%    & 10\%    & 24\% \\
Single fix commit      & 91\%    & 83\%    & 80\%    & 96\% \\
Gold files (median)    & 1       & 1       & 2       & 1 \\
\midrule
Memory pool            & 4{,}033 & 6{,}014 & 19{,}820 & 2{,}315 \\
\bottomrule
\end{tabular}
\end{table}

\Cref{tab:datasheet} reports the yield of every condition. The final
benchmark holds 6{,}437 instances, of which 6{,}314 are scored; the
123 we leave unscored, and why, are described in
\cref{sec:setup:populations}. The largest loss differs by project
and is informative in itself. Qt loses most instances at the linkage
step, because its tracker covers about forty repositories and we mine
only \texttt{qtbase}. MongoDB loses most at the source-file step, because
many of its fixes change only test or build files.

\subsection{Easy and hard instances}
\label{sec:benchmark:strata}

The challenge notes that error logs often show where a failure was
detected rather than where it originated. Our benchmark makes this
measurable. We call an instance \emph{easy} if the issue text mentions
the name of at least one gold file, for example in a pasted stack trace.
We call it \emph{hard} otherwise: the report describes the failure, but
no gold file is named anywhere in it. Across the benchmark, 53\% of
instances are hard. The fraction varies strongly by project
(\cref{tab:datasheet}). Kafka reports often contain a Java stack trace
that names the faulty class, so only 28\% are hard. Qt reports describe
misbehaving user interfaces and rarely contain a trace, so 68\% are hard.
Every result in this paper is reported on both strata. We also record
whether each instance contains a stack trace at all, which supports the
same analysis from the input side.

The definition is file-based, and that choice is not language-neutral:
a Java frame names the file (class and file coincide), while a C++
frame names a function. To bound this artifact we computed a second
definition that also counts a mention of a gold file's basename stem or
of any function defined in a gold file. The hard fraction then falls to
14\% (Kafka), 33\% (MariaDB), 38\% (MongoDB), and 38\% (Qt): part of
the Java-to-C++ gap is indeed the naming convention, and the three C++
projects converge to a band of 33\% to 38\% while Kafka stays well
below it. We
keep the file-based definition for stratification because a file name
is what grep turns into a candidate file directly, but both labels ship
with the benchmark.

\subsection{The incident memory pool}
\label{sec:benchmark:memory}

Our approach uses the project's resolved history as a signal
(\cref{sec:approach}). This history does not need logs: any fixed issue
whose fix commit we can identify tells us which files a past problem
ended in. We therefore build a separate \emph{memory pool} per project
from all of PJD, using conditions 2 to 4 but not condition 1. Each pool
entry carries the issue's summary and description, its creation and
resolution timestamps, its fix commits, and the gold files of those
commits. The pools are between three and six times larger than the benchmarks
(\cref{tab:datasheet}). PJD's typed issue links also give us duplicate
clusters: groups of issues marked as duplicates of one another. We use
the timestamps and the clusters to keep the memory honest during
evaluation (\cref{sec:approach:firewall}).

\subsection{Limitations of mined ground truth}
\label{sec:benchmark:limits}

Mined gold is not perfect, and we state its known weaknesses. A commit
subject can mention an issue it does not fix, although our filters remove
the most common cases (reverts, merges, and multi-commit efforts). A fix
can touch files beyond the faulty one, for example during a small
refactoring; the median fix modifies one to two source files, which
bounds this effect. Finally, the file a fix modifies is the place the
developers chose to repair, which is occasionally not the place the
defect originated. We accept these imperfections because they are the
same imperfections real resolved histories have. A method that works on
this benchmark works on what projects actually record.
